% ============================================================
% CAPITOLO 26 - SETTIME DI SECONDA SPECIE
% ============================================================
\chapter{Settime di seconda specie}
\label{chap:seconda_specie}

La seconda specie è formata da una \textbf{triade minore} e una \textbf{settima minore}. L'accordo di gran lunga più importante di questa specie è quello costruito sul \textbf{secondo grado} del modo maggiore.

\section{Funzione e risoluzione}
\label{sec:seconda_specie_dettaglio}

La funzione tipica è di \textbf{sottodominante}. Poiché non è un accordo di dominante, la settima deve essere rigorosamente \textbf{preparata}. L'accordo risolve naturalmente sulla dominante (V o \sigla{V}{7}{}).

\begin{regola}[Collegamento I-\sigla{II}{7}{}-V e raddoppio della terza]
    \label{reg:ii7_raddoppio}
    Nel collegamento tra il I grado e il \sigla{II}{7}{}, e nella successiva risoluzione sul V, è frequente e consigliabile \textbf{raddoppiare la terza} dell'accordo di \sigla{II}{7}{} (che corrisponde al IV grado della scala). Serve a evitare \textbf{quinte parallele} tra la fondamentale e la quinta del II grado nel movimento verso la dominante.

	In alternativa, o in aggiunta, si può \textbf{omettere la quinta} del \sigla{II}{7}{}, raddoppiando un'altra nota e lasciando di fatto la terza raddoppiata.
\end{regola}

\begin{esempio}[\sigla{II}{7}{} in Do maggiore]
    \musica{es_cap13_minore_do_1}
    \musica{es_cap13_minore_do_2}
    Il \sigla{II}{7}{} risolve sul V: la settima Do scende a Si (sensibile), che non si raddoppia. Il primo rivolto (\sigla{II}{6}{5}, basso Fa) sale a Sol, il secondo (\sigla{II}{4}{3}, basso La) scende a Sol, il terzo (\sigla{II}{2}{}, basso Do) scende a Si sul \sigla{V}{6}{}.
\end{esempio}

\begin{esempio}[Seconda specie in La minore: \sigla{I}{7}{} $\rightarrow$ IV]
	\label{ex:i7_minore}
	Nel minore naturale la settima minore sul I grado (La-Do-Mi-Sol) risolve sul IV (Re-Fa-La): la settima Sol scende a Fa, terza del IV. Il comportamento è lo stesso dell'esempio in Do maggiore.
	\musica{es_cap13_minore_la_1}
	\musica{es_cap13_minore_la_2}
\end{esempio}

\section{Collocazione sulle scale}
\label{subsec:ii7_collocazione}

\begin{regola}[Gradi della seconda specie]
	\label{reg:ii7_gradi}
	L'accordo di settima minore (seconda specie) si trova:
	\begin{itemize}
		\item sul \textbf{II}, \textbf{III} e \textbf{VI} grado della scala maggiore;
		\item sul \textbf{IV} grado della scala minore armonica (nel minore naturale anche sul I e sul V).
	\end{itemize}
	Il più importante è quello sul \textbf{II grado del modo maggiore}, con funzione di sottodominante.
\end{regola}

\section{Stato fondamentale}
\label{subsec:ii7_stato_fondamentale}

\begin{regola}[\sigla{II}{7}{} in stato fondamentale]
	\label{reg:ii7_stato_fondamentale}
	Con il II grado al basso la settima va \textbf{preparata} nell'accordo precedente (presente nella stessa voce, con durata pari o superiore) e \textbf{tenuta}, poi risolve \textbf{scendendo per grado congiunto}. Le altre note, compresa la fondamentale, sono \textbf{parti libere} e si possono raddoppiare. La quinta si può omettere raddoppiando un'altra nota (mai la settima).
\end{regola}

Il \sigla{II}{7}{} risolve naturalmente sul \textbf{V grado}, armonizzato con la triade semplice, con la cadenza composta consonante o dissonante, oppure con la \sigla{V}{7}{}.

\section{I tre rivolti}
\label{subsec:ii7_rivolti}

\begin{center}
\begin{tabularx}{\linewidth}{lLLLL}
\toprule
\textbf{Stato} & \textbf{Cifratura} & \textbf{Grado al basso} & \textbf{Risolve su} & \textbf{Basso obbligato} \\
\midrule
Fond. & 7 & II & V (o \sigla{V}{7}{}) & no \\
I rv. & 6/5 & IV & V & sì (sale) \\
II rv. & 4/3 & VI & V & sì (scende) \\
III rv. & 4/2 & I (settima) & \sigla{V}{6}{} o \sigla{V}{6}{5} & sì (scende) \\
\bottomrule
\end{tabularx}
\end{center}

\begin{itemize}
	\item \textbf{\sigla{II}{6}{5}} (IV grado al basso): il basso risolve per grado congiunto ascendente sulla dominante. Si impiega molto nella formula cadenzale I (oppure \sigla{I}{6}{} o VI) $-$ \sigla{II}{6}{5} $-$ \sigla{I}{6}{4} $-$ V. Può risolvere anche sulla triade o sulla settima di sensibile in stato di rivolto.
	\item \textbf{\sigla{II}{4}{3}} (VI grado al basso): il basso risolve per moto discendente sul V grado.
	\item \textbf{\sigla{II}{2}{}} (I grado al basso, la settima): la settima al basso va preparata e risolve per moto discendente sul VII grado, armonizzato con il primo rivolto della dominante (\sigla{V}{6}{} o \sigla{V}{6}{5}).
\end{itemize}

\begin{osservazione}[Progressioni e IV grado del modo minore]
	\label{oss:ii7_progressioni}
	\begin{itemize}
		\item Il \sigla{VI}{7}{} e le altre settime di seconda specie si usano scolasticamente nelle \textbf{progressioni discendenti} (per esempio $\sigla{III}{7}{} - \sigla{VI}{7}{} - \sigla{II}{7}{}$).
		\item La seconda specie sul IV grado del modo minore è un accordo \textbf{derivato} (da una nona di terza specie sul II grado senza fondamentale) e si comporta come la quarta specie sul IV grado del maggiore.
	\end{itemize}
\end{osservazione}
